\chapter{Detalle de la validación externa}
\label{anx:validacion}

\section{Carga aislada}

Cada conjunto externo se carga con su identificador de ejecución, y de ahí salen sus propios
\textit{topics} (\texttt{raw.fa.<id>.*}, \texttt{norm.fa.<id>.*} y los de detección) y sus grupos
de consumo. Los pasos son siempre los mismos:
\begin{enumerate}
    \item foto del almacén antes de cargar;
    \item reproducción y normalización, comprobando que lo crudo es igual a lo normalizado más lo
    rechazado con motivo;
    \item detección por tandas;
    \item volcado a ClickHouse;
    \item patrones en Flink con el plan sin verdad;
    \item foto de nuevo.
\end{enumerate}
La carga solo se da por buena si las dos fotos de los 110 conjuntos del Pilar~A coinciden. La
segunda versión de los patrones lee los mismos \textit{topics} normalizados y escribe en
\texttt{det.v2.<grupo>.*} y en una tabla de campañas propia, así que v1 y v2 se comparan sobre los
mismos eventos.

\section{COMISET: censo y selección}

El archivo de la red real ocupa 31,7\,GB comprimido y 934,5\,GB descomprimido. Se recorrió en
flujo, sin descomprimirlo a disco, y el censo completo tardó 59 minutos: 202,3 millones de registros,
16 equipos y 14 días, de los que el 99,98\,\% son de Sysmon. Antes de detectar se fijó la
selección:
\begin{itemize}
    \item 8 horas por equipo, repartidas por toda su actividad;
    \item un tope de 20.000 registros por hora.
\end{itemize}
Salen 110 horas y 1,75 millones de registros, y cada hora es un conjunto. El 99,91\,\% se normaliza;
los rechazos son tipos de Sysmon sin mapeador (DNS, tuberías y otros).

En el laboratorio de COMISET, el criterio también se fijó antes de detectar:
\begin{itemize}
    \item las horas de ataque son las de más etiquetas de las técnicas que miran los patrones;
    \item las de control, horas sin ninguna de esas etiquetas;
    \item cada hora de ataque se recorta a la actividad etiquetada, con 10 minutos de margen.
\end{itemize}
Las etiquetas resultaron ser sobre todo ruido de las reglas que las ponen: el 95\,\% de las de
inyección son el Explorador abriendo el Administrador de tareas, las de credenciales del registro son
escrituras de Restaurar sistema y los accesos a LSASS son de servicios del sistema con permisos
inocuos. Elegir por número de etiquetas eligió horas ruidosas, no horas de ataque. El ataque real de
la muestra (la explotación de un servicio de control remoto) da dos alertas de nivel alto; la
elevación no la cubre ninguna regla, y se deja sin corregir porque el conjunto es de validación. El
presupuesto de unos 600.000 eventos lo impuso la detección. Un solo equipo cae en una única
partición y se detecta con un solo proceso: unos 255 eventos por segundo con v1 y unos 90 con v2.
Hizo falta poder reanudar la detección sin duplicar.

\section{Laboratorio propio}

La telemetría del laboratorio no trae número de registro, y su hora está en la zona local de la
máquina virtual. El adaptador hace dos cosas:
\begin{itemize}
    \item convierte la hora a UTC con la zona declarada y la comprueba contra el campo de hora UTC
    de Sysmon, con un desfase exacto de 120 minutos en los 11.959 eventos que lo traen;
    \item toma como identidad del evento la línea del fichero de origen.
\end{itemize}
Se normaliza el 99,5\,\%; solo quedan fuera las 81 consultas DNS de Sysmon. El corte entre el ataque
y la línea base se fijó 37 minutos después de la última ventana de ataque.

La máquina virtual se borró durante el desarrollo, pero los datos se conservaban. El guion de
etiquetado se reconstruyó a partir de ellos, y reproduce fila a fila el conjunto final de 14.895
eventos. El guion que orquestaba la captura se perdió, así que la captura no puede repetirse
exactamente igual.

\begin{table}[H]
\centering
\small
\begin{tabular}{|l|r|l|}
\hline
\textbf{Técnica (pruebas)} & \textbf{Alertas en ventana} & \textbf{Reglas específicas} \\
\hline
\texttt{T1547.001} (19) & 138 & Carpeta de inicio, claves de ejecución automática \\
\texttt{T1070} (2) & 12 & \texttt{fsutil} \\
\texttt{T1055} (3) & 15 & Ninguna (solo el andamiaje del marco de pruebas) \\
\texttt{T1003.001} (5) & 40 & \texttt{procdump} sobre LSASS, \texttt{comsvcs.dll} \\
\hline
\end{tabular}
\caption{Alertas dentro de las ventanas del manifiesto. Parte de ellas son el reconocimiento con
\texttt{whoami} y \texttt{hostname} que lanza el propio marco de pruebas.}
\end{table}

La campaña de P1 empieza 13 minutos antes del volcado. La primera evidencia es una copia de
seguridad del registro hecha por una prueba de persistencia, que una regla de volcado de colmenas
del registro confunde con un volcado. Es imprecisión de la regla, no del patrón. En la línea base,
518 de las 553 alertas son reconocimiento de bajo nivel (\texttt{whoami} y \texttt{systeminfo}) que
lanza el generador de actividad simulada, y ninguna abre campaña. La línea base incluye además un
reinicio de la máquina en el que se ejecuta la persistencia que dejó el ataque: sus seis alertas de
nivel alto son verdaderos positivos, no falsas alarmas.
